\begin{proposition}[geometry of the family]
\label{prop:family-geometry}
The projection $g$ is proper, $\mathcal X$ is a smooth complex threefold,
  and $S$ is a reduced irreducible normal-crossing (self-crossing) surface whose normalization
is $\P^1\times\P^1$.  Its singular locus is
\(D=\{[0:0:1]\}\times\P^1\), and locally along $D$ the pair $(\mathcal X,S)$
has the model (3.1); away from $D$ the fibre is smooth and has the usual
local model $(\C^3,\{u=0\})$.
\end{proposition}

\begin{proof}
Properness follows because $\mathscr C$ is closed in the projective
projection $\P^2\times\Delta\to\Delta$.  Write the defining equation as
\[
 F=y^2z-x^2(x+z)-tz^3.
\]
On the chart $z\ne0$, $\partial F/\partial t=-z^3\ne0$, so $\mathscr C$
is smooth there.  On $z=0$, the equation forces $x=0$ and $y\ne0$; at such
points $\partial F/\partial z=y^2\ne0$.  Hence $\mathscr C$, and therefore
$\mathcal X$, is smooth.  The only finite discriminant values of the fibres
are $t=0$ and $t=-4/27$; after shrinking $\Delta$ to $|t|<4/27$, every
noncentral fibre is smooth.

The central curve is
\(C_0:\ y^2z=x^2(x+z)\).  Its only singular point is
\(n=[0:0:1]\): on $z=1$ the equations $F=F_x=F_y=0$ give $(x,y)=(0,0)$,
and the point at infinity $[0:1:0]$ is smooth because $F_z=1$ there
(Euler's identity also gives the remaining projective singularity equation).
The parametrization
\[
 \nu:\P^1\longrightarrow C_0,\qquad
 [u:v]\longmapsto
 [(u^2-v^2)v:\,u(u^2-v^2):\,v^3]
 \tag{4.3}
\]
maps the two points $p=[1:1]$, $q=[-1:1]$ to $n$.  On the chart $v\ne0$, set
$r=u/v$; after scaling, its image is
$[r^2-1:r(r^2-1):1]$.  Conversely, for an affine point of $C_0$ with
$x\ne0$, the equation gives $r=y/x$ and $x=r^2-1$, so it is hit by
$[r:1]$.  The only affine point with $x=0$ is the node $(0,0)$, hit by
$[1:1]$ and $[-1:1]$, and the point at infinity $[0:1:0]$ is hit by
$[1:0]$.  Thus $\nu(\P^1)=C_0$; as a nonconstant morphism to a projective
curve it is finite and generically one-to-one, hence the normalization.  In
particular $C_0$ is irreducible and rational.  The product normalization
is $\eta=\nu\times\id_{\P^1}:\P^1\times\P^1\to S$.

At $n$ the affine equation is
\(t=y^2-x^2-x^3\).  On a sufficiently small neighbourhood choose the
holomorphic square root $h(x)=\sqrt{1+x}$ and set
\(u=y-xh(x)\), \(v=y+xh(x)\).  Then $uv=t$, and the map
$(x,y)\mapsto(u,v)$ has invertible $2\times2$ Jacobian at $(0,0)$, so $u,v$
are local coordinates on the surface $\mathscr C$.  If $s$ is a local
coordinate on $\P^1$, then $(u,v,s)$ are coordinates on $\mathcal X$ and
$t=uv$.  Hence the pair has the asserted normal-crossing (self-crossing)
local model, with singular locus $D$ as stated.
\end{proof}

